\section{国内外研究现状及分析}

\begin{frame}{国内外研究现状及分析}
    \textbf{RS码理论研究综述：早期奠定基础的编解码理论}
    
    \vspace{0.5em}

    \begin{enumerate}
        \item 由Reed和Solomon提出\cite{reedPolynomialCodesCertain1960}，基于有限域Vandermonde矩阵乘法。
        \item Luby等人将乘法变换为异或，并使用求逆更快的Cauchy矩阵\cite{lubyXORBasedErasureResilientCoding1995}。
    \end{enumerate}
    
    这种理论最优的RS码是后续优化研究的基础。

    \vspace{0.5em}

    \textbf{RS码优化研究综述：重点研究工作及其优化方案总结}
    
    \vspace{0.5em}

    \begin{table}[]
    \centering
    \begin{tabular}{|c|c|}
    \hline
    减少矩阵中$1$的个数进而减少运算 & Zerasure\cite{zhouFastErasureCoding2020}、Cerasure\cite{wangFastAccelerationStrategies2025} \\ \hline
    提取并复用公共异或运算表达式    & SLPEC\cite{uezatoAcceleratingXORbasedErasure2021}、Cerasure                                 \\ \hline
    内存与缓存访问优化  & Zerasure、SLPEC、Cerasure \\ \hline
    降低数据指针相关开销 & Cerasure                \\ \hline
    \end{tabular}
    \end{table}

    Cerasure综合各优化策略，并提出指针开销优化，是RS码SOTA实现。

\end{frame}

\begin{frame}{国内外研究现状及分析}

    \textbf{移位异或纠删码理论研究综述：从ZigZag到Two-tone}
    
    \vspace{1em}

    \begin{columns}[t]
        \column{0.5\textwidth}
        ZigZag Decodable Codes\cite{gongZigzagDecodableCodes2017}是较早提出的移位异或纠删码，通过构建特殊的移位矩阵实现了仅依赖异或和移位的高效编解码
        \column{0.5\textwidth}
        twotone
    \end{columns}
    
\end{frame}
